\documentclass[11pt,a4paper]{article}

% ---------------------------------------------------------
% PREAMBLE (stable, warning-free, Overleaf-safe)
% ---------------------------------------------------------
\usepackage[utf8]{inputenc}
\usepackage[T1]{fontenc}
\usepackage{lmodern}
\usepackage{amsmath, amssymb, amsfonts}
\usepackage{bm}
\usepackage{geometry}
\usepackage{graphicx}
\usepackage{hyperref}
\usepackage{cite}
\usepackage{authblk}
\usepackage{physics}
\usepackage{mathtools}

\geometry{margin=2.4cm}

\title{\textbf{Companion XLI: The Epoch of Reionization in the Relaxation-Driven Cyclic Cosmology}}
\author{Michael Lehmann \\ \texttt{mi.lehmann@gmx.de}}
\date{April 2026}

\begin{document}
\maketitle

\begin{abstract}
The Epoch of Reionization (EoR) marks the transition from a neutral to an ionized 
Universe and represents one of the most sensitive probes of early structure 
formation. In the standard cosmological model, reionization is driven by the 
formation of the first stars and galaxies, whose ultraviolet radiation ionizes the 
intergalactic medium. In the Relaxation-Driven Cyclic Cosmology (RDCC), the 
scale-dependent evolution of the gravitational potential modifies the formation 
history of early galaxies, the propagation of ionization fronts, and the topology 
of reionization. This Companion develops a continuous description of reionization 
in RDCC, deriving how the relaxation scale influences the timing, morphology, and 
observational signatures of the EoR. The resulting framework provides a natural 
interpretation of the early appearance of massive galaxies and the extended 
ionization structures observed in deep surveys.
\end{abstract}
% ---------------------------------------------------------
\section{Introduction}

The Epoch of Reionization is a pivotal phase in cosmic history. After recombination, 
the Universe remained neutral for several hundred million years, until the first 
stars and galaxies formed and emitted ultraviolet photons capable of ionizing 
hydrogen. The resulting ionized bubbles expanded, merged, and eventually filled the 
Universe.

In the standard cosmological model, the timing and morphology of reionization are 
determined by the formation rate of early galaxies and the clumpiness of the 
intergalactic medium. In RDCC, the gravitational potential evolves differently. The 
relaxation mechanism suppresses the growth of small-scale modes and enhances the 
coherence of large-scale modes. This modifies the formation history of early 
galaxies and the propagation of ionization fronts, imprinting the relaxation scale 
$k_{\mathrm{rel}}$ onto the structure of the EoR.
% ---------------------------------------------------------
\section{Early Galaxy Formation in RDCC}

The formation of the first galaxies is governed by the collapse of dark matter 
halos. In RDCC, the gravitational potential takes the form
\begin{equation}
    \Phi_{\mathrm{RDCC}}(k,a)
    = \Phi(k,a)
      \left[
        1 - \chi_{\Phi}
        \left(\frac{k}{k_{\mathrm{rel}}}\right)^{\alpha}
      \right].
\end{equation}

Large-scale modes, which dominate the collapse of massive halos, retain their 
coherence. As a result, massive halos form earlier in RDCC than in the standard 
cosmological model. Small-scale modes, which would normally produce numerous 
low-mass halos, are suppressed, delaying the formation of dwarf galaxies.

The halo mass function at high redshift becomes
\begin{equation}
    n_{\mathrm{RDCC}}(M,z)
    = n(M,z)
      \left[
        1 - \chi_{n}
        \left(\frac{k(M)}{k_{\mathrm{rel}}}\right)^{\alpha}
      \right].
\end{equation}

This modification naturally explains the early appearance of massive galaxies 
observed by JWST.
% ---------------------------------------------------------
\section{Ionization Fronts and the RDCC Radiation Field}

Ionization fronts propagate through the intergalactic medium according to
\begin{equation}
    \frac{dR_{\mathrm{I}}}{dt}
    = \frac{\dot{N}_{\gamma}}{4\pi R_{\mathrm{I}}^{2} n_{\mathrm{H}}}
      - \alpha_{\mathrm{B}} n_{\mathrm{H}} R_{\mathrm{I}},
\end{equation}
where $\dot{N}_{\gamma}$ is the ionizing photon production rate.

In RDCC, the altered formation history of early galaxies modifies $\dot{N}_{\gamma}$. 
Massive galaxies, which form earlier, produce strong ionizing radiation fields, 
while the suppressed formation of dwarf galaxies reduces the contribution from 
low-mass halos.

The resulting ionization fronts are larger, smoother, and more coherent than in the 
standard cosmological model. Their expansion reflects the coherence of the 
gravitational potential on scales near $k_{\mathrm{rel}}$.
% ---------------------------------------------------------
\section{Topology of Reionization in RDCC}

The topology of reionization describes the spatial distribution of ionized and 
neutral regions. In RDCC, the suppression of small-scale modes reduces the 
fragmentation of ionized bubbles, while the enhanced coherence of large-scale modes 
promotes the formation of extended ionized regions.

The ionized volume fraction evolves as
\begin{equation}
    Q_{\mathrm{HII,RDCC}}(z)
    = Q_{\mathrm{HII}}(z)
      \left[
        1 + \chi_{Q}
        \left(\frac{k_{\mathrm{rel}}}{k_{\mathrm{HII}}}\right)^{\alpha}
      \right],
\end{equation}
where $k_{\mathrm{HII}}$ characterizes the typical size of ionized bubbles.

This modification produces a smoother reionization topology with larger ionized 
regions and fewer small-scale neutral pockets.
% ---------------------------------------------------------
\section{21-cm Signatures and Observational Probes}

The 21-cm line of neutral hydrogen provides a direct probe of the EoR. The 
brightness temperature contrast is given by
\begin{equation}
    \delta T_{b}
    \propto x_{\mathrm{HI}}
    \left(1 - \frac{T_{\gamma}}{T_{s}}\right),
\end{equation}
where $x_{\mathrm{HI}}$ is the neutral fraction and $T_{s}$ is the spin temperature.

In RDCC, the modified reionization topology and the altered formation history of 
early galaxies imprint distinctive signatures on the 21-cm power spectrum. The 
suppression of small-scale modes reduces small-scale fluctuations, while the 
enhanced coherence of large-scale modes amplifies large-scale structure.

These signatures provide a direct observational window into the relaxation scale 
and the temporal structure of the gravitational field.
% ---------------------------------------------------------
\section{Conclusion}

The Epoch of Reionization in the Relaxation-Driven Cyclic Cosmology reflects the 
scale-dependent evolution of the gravitational potential. The relaxation mechanism 
modifies the formation history of early galaxies, the propagation of ionization 
fronts, and the topology of reionization. These effects produce distinctive 
signatures in the ionized volume fraction, the morphology of ionized bubbles, and 
the 21-cm power spectrum. The present Companion establishes the theoretical 
foundation for future studies of early galaxy formation, cosmic dawn, and the 
interaction between radiation and matter in RDCC.

% ========================
% References
% ========================

\begin{thebibliography}{10}

\bibitem{Flagship} Lehmann, M. (2026). Relaxation-Driven Cyclic Cosmology (RDCC): A Minimal CPT-Symmetric Framework with a Single Parameter. Flagship Paper. Zenodo: \url{https://zenodo.org/records/19744958} (v43.0 or later).

\bibitem{Zenodo} The complete RDCC companion series (I--LVII) is available at the same Zenodo record above.

% Thematisch relevante Companions
\bibitem{CompXVI} Companion XVI: Boltzmann Code and Modified Hierarchy.
\bibitem{CompXXXI} Companion XXXI: RDCC Halo Model and Non-Linear Structure.
\bibitem{CompXXXII} Companion XXXII--XXXIX: Galaxy--Halo Connection, Cosmic Web, Lensing, RSD, ISW, tSZ, Rotation Curves, CGM.

% Wichtige externe Referenzen
\bibitem{Planck2020} Planck Collaboration (2020), Astron. Astrophys. 641, A6.
\bibitem{DESI2024} DESI Collaboration (2024/2025), arXiv: [aktuelle $f\sigma_8$/BAO-Paper].
\bibitem{JWST2024} Maiolino et al. (2024), Nature (JWST high-$z$ galaxies/quasars).
\bibitem{Kaiser1987} Kaiser, N. (1987), Mon. Not. R. Astron. Soc. 227, 1 (RSD).
\bibitem{Bartelmann2001} Bartelmann, M. \& Schneider, P. (2001), Phys. Rep. 340, 291 (Weak Lensing Review).
\bibitem{HuOkamoto2002} Hu, W. \& Okamoto, T. (2002), Astrophys. J. 574, 566 (CMB Lensing).
\bibitem{LewisChallinor2006} Lewis, A. \& Challinor, A. (2006), Phys. Rep. 429, 1 (CMB Lensing Review).
\bibitem{NFW1997} Navarro, J. F., Frenk, C. S. \& White, S. D. M. (1997), Astrophys. J. 490, 493 (Halo Profiles).

\end{thebibliography}

\end{document}
